\documentclass[
10pt, % Main document font size
a4paper, % Paper type, use 'letterpaper' for US Letter paper
oneside, % One page layout (no page indentation)
%twoside, % Two page layout (page indentation for binding and different headers)
headinclude,footinclude, % Extra spacing for the header and footer
BCOR5mm, % Binding correction
]{scrartcl}
%%%%%%%%%%%%%%%%%%%%%%%%%%%%%%%%%%%%%%%%%
% Arsclassica Article
% Structure Specification File
%
% This file has been downloaded from:
% http://www.LaTeXTemplates.com
%
% Original author:
% Lorenzo Pantieri (http://www.lorenzopantieri.net) with extensive modifications by:
% Vel (vel@latextemplates.com)
%
% License:
% CC BY-NC-SA 3.0 (http://creativecommons.org/licenses/by-nc-sa/3.0/)
%
%%%%%%%%%%%%%%%%%%%%%%%%%%%%%%%%%%%%%%%%%

%----------------------------------------------------------------------------------------
%	REQUIRED PACKAGES
%----------------------------------------------------------------------------------------

\usepackage[
nochapters, % Turn off chapters since this is an article        
beramono, % Use the Bera Mono font for monospaced text (\texttt)
eulermath,% Use the Euler font for mathematics
pdfspacing, % Makes use of pdftex’ letter spacing capabilities via the microtype package
dottedtoc % Dotted lines leading to the page numbers in the table of contents
]{classicthesis} % The layout is based on the Classic Thesis style

\usepackage{arsclassica} % Modifies the Classic Thesis package
\usepackage{ulem}
\usepackage{graphicx}
\usepackage{float}
\usepackage{algorithm2e}
\usepackage[export]{adjustbox}
\usepackage[utf8]{inputenc}  
\usepackage[english,frenchb]{babel} 
\usepackage[T1]{fontenc} % Required for accented characters
\usepackage{multicol}
\usepackage{graphicx} % Required for including images
\graphicspath{{Figures/}} % Set the default folder for images

\usepackage{enumitem} % Required for manipulating the whitespace between and within lists

\usepackage{lipsum} % Used for inserting dummy 'Lorem ipsum' text into the template


\expandafter\def\csname ver@subfig.sty\endcsname{}


\usepackage{svg}
\usepackage{hyperref}

\usepackage{amsmath,amssymb,amsthm} % For including math equations, theorems, symbols, etc

\usepackage{varioref} % More descriptive referencing

%----------------------------------------------------------------------------------------
%	THEOREM STYLES
%---------------------------------------------------------------------------------------

\theoremstyle{definition} % Define theorem styles here based on the definition style (used for definitions and examples)
\newtheorem{definition}{Definition}

\newtheorem*{remark}{Remark}
\theoremstyle{plain} % Define theorem styles here based on the plain style (used for theorems, lemmas, propositions)
\newtheorem{theorem}{Theorème}
\newtheorem{lemma}[theorem]{Lemme}
\newtheorem{proposition}[theorem]{Proposition}

\theoremstyle{remark} % Define theorem styles here based on the remark style (used for remarks and notes)

\DeclareMathOperator*{\argmin}{arg\,min}
\DeclareMathOperator*{\argmax}{arg\,max}

%----------------------------------------------------------------------------------------
%	HYPERLINKS
%---------------------------------------------------------------------------------------

\hypersetup{
%draft, % Uncomment to remove all links (useful for printing in black and white)
colorlinks=true, breaklinks=true, bookmarks=true,bookmarksnumbered,
urlcolor=webbrown, linkcolor=RoyalBlue, citecolor=webgreen, % Link colors
pdftitle={}, % PDF title
pdfauthor={\textcopyright}, % PDF Author
pdfsubject={}, % PDF Subject
pdfkeywords={}, % PDF Keywords
pdfcreator={pdfLaTeX}, % PDF Creator
pdfproducer={LaTeX with hyperref and ClassicThesis} % PDF producer
}

\hyphenation{Fortran hy-phen-ation} 

\title{\normalfont{FTML Exercices 4}\\
} 
\author{\spacedlowsmallcaps{}
} 
\date{\today}
\begin{document}
\renewcommand{\sectionmark}[1]{\markright{\spacedlowsmallcaps{#1}}}
\lehead{\mbox{\llap{\small\thepage\kern1em\color{halfgray} \vline}\color{halfgray}\hspace{0.5em}\rightmark\hfil}} 
\pagestyle{scrheadings}
\maketitle 
\setcounter{tocdepth}{3} 

\tableofcontents

\section{\large\color{Blue}Forte convexité de la loss Ridge regression}

La forte convexité est une version plus forte de la convexité: la fonction considérée doit non seulement etre convexe, mais également avoir une certaine courbure minimale. Cette propriété est utile par exemple pour avoir des résultats de vitesse de convergence des descentes de gradient sur des fonctions fortement convexes.

On rappelle la définition de la convexité et de la forte convexité:

\begin{definition}{Convex function}
    \label{def:convexity}

    The function $f:\Omega \rightarrow \mathbb{R} $ with $\Omega$ a convex subset of $V$ is:
    \begin{itemize}
	\item \textbf{{convex}} if $\forall x, y\in \Omega, \alpha\in [0,1]$,
	    \begin{equation*}
		f(\alpha x+(1-\alpha)y)\leq \alpha f(x)+(1-\alpha)f(y)
	    \end{equation*}
	\item \textbf{{strictly convex}} if $\forall x, y\in \Omega, \alpha\in [0,1]$,
	    \begin{equation*}
		f(\alpha x+(1-\alpha)y)< \alpha f(x)+(1-\alpha)f(y)
	    \end{equation*}
	\item $\mu$-\textbf{{strongly convex}} if $\forall x, y\in \Omega, \alpha\in [0,1]$,
	    \begin{equation*}
		f(\alpha x+(1-\alpha)y)\leq \alpha f(x)+(1-\alpha)f(y)- \frac{\mu}{2} \alpha(1-\alpha)||x-y||^2
	    \end{equation*}
    \end{itemize}
\end{definition}

On souhaite montrer que la loss Ridge regression est fortement convexe en son argument $\theta$.

\begin{definition}{Ridge regression estimator}

\begin{equation}
    \label{risk:ridge}
    \hat{ \theta}_{\lambda} = \argmin_{\theta \in \mathbb{R}^d} \Big( \frac{1}{n} ||y-X\theta||_2^2 + \lambda || \theta||_2^2 \Big)
\end{equation}
\end{definition}


        \begin{itemize}
            \item 1) Pour s'entrainer à manipuler la forte convexité, vérifier que pour $x\in \mathbb{R}$, l'application $x\mapsto x^2$ est $2$-fortement convexe.
            \item 2) Montrer que pour $x\in \mathbb{R}^d$, l'application $x\mapsto ||\theta||^2$ est $2$-fortement convexe.
            \item 3) En déduire le résultat sur la forte convexité de la loss Ridge.
            \item 4) (Un peu plus difficile, mais pas nécessaire pour avoir le résultat sur Ridge) Prouver le résultat plus général suivant:

     if $Q$ is a symmetric definite positive matrix (matrice définie positive) with smallest eigenvalue $\lambda_{min}>0$, then the positive definite quadratic form $x\mapsto x^TQx$ is $ 2\lambda_{min}$- strongly convex.
        \end{itemize}


\end{document}
